The counting script enumerates every retained triple $(a,l_2,l_3)$ with
$a+2l_2+2l_3\le 2N$, using the preceding character dimensions.  The following
counts apply to the current physical PBW reference code, before any numerical
values are computed:
\begin{center}
\begin{tabular}{lrr}
\toprule
Quantity & $N=10$ & $N=15$\\
\midrule
Physical monomials & 506 & 1,496\\
Parity output slots & 4,048 & 11,968\\
Largest NS Gram dimension & 161 & 1,016\\
Largest R Gram dimension, per parity & 232 & 1,472\\
Entries in the distinct Gram matrices & 450,469 & 20,431,839\\
Sum of cubes of Gram dimensions & 78,765,711 & 22,116,449,163\\
Requested vertex entries, $\eta\eta'=-1$ & 315,028 & 9,895,892\\
Dense contraction multiplications & 7,204,872 & 734,995,948\\
\bottomrule
\end{tabular}
\end{center}
For a triple with dimensions $d_1,d_2,d_3$, the two parity sectors and two
vertices request $4d_1d_2d_3$ entries; the three successive Gram contractions
and final pairing use $2d_1d_2d_3(d_1+d_2+d_3+1)$ complex multiplications.
The Gram cache has one matrix at each NS half-level and two parity matrices
at each level on each Ramond edge.  The cubic sum is a work proxy for dense
inversion, not an exact operation count.  At level 15, storing only these
Gram matrices and their inverses as complex128 arrays already takes
$654$ MB; Python objects, Ward caches, and multiprecision storage are additional.

For $\eta'=\eta$, both vertices call the same cached three-point form.
The number of distinct requested form entries is therefore halved.  The
implementation still constructs both tensor arrays, including their phases,
and performs the same Gram contractions.  Consequently the whole PBW runtime
must not be halved.  In the estimates below only the fitted cost associated
with tensor entries is allowed to decrease, by a factor between $1/2$ and $1$.
The endpoints describe assumptions about the unmeasured division between Ward
work and tensor/cache overhead; they are not rigorous bounds.

The calibration uses the 506 cumulative records in each of
\path{Code/theta_fermion_ccy/results/pbw_machine_level10.log} and
\path{Code/theta_fermion_ccy/results/pbw_level10.log}.
Differences of successive records give times per monomial.  A nonnegative
least-squares fit uses newly encountered Gram entries, requested tensor
entries, newly encountered sums of dimension cubes, dense contraction
multiplications, monomial count, and cumulative output size.  Evaluating the
same features at level 15 gives the following \emph{full-feature estimates};
no PBW calculation was rerun. A leaner model below gives substantially lower
level-15 estimates, so the table is not a definitive runtime forecast.
\begin{center}
\begin{tabular}{clcc}
\toprule
Level & PBW arithmetic & $\eta\eta'=1$ & $\eta\eta'=-1$\\
\midrule
10 & Native, 53 bits & $21$--$31$ s (estimated) & $31.4$ s (saved)\\
15 & Native, 53 bits & $34$--$40$ min (estimated) & $40$ min (estimated)\\
10 & FLINT, 384 bits & $89$--$110$ s (estimated) & $110.2$ s (saved)\\
15 & FLINT, 384 bits & $241$--$252$ min (estimated) & $252$ min (estimated)\\
\bottomrule
\end{tabular}
\end{center}
The level-10 saved full-process times are $33.55$ s and $112.83$ s,
respectively.  Their internal times include repeatedly serializing the growing
coefficient list; the C++ internal timer excludes final serialization, so the
parent process times are the appropriate end-to-end comparison.

The time model is substantially less certain than the integer work counts.
A lean fit omitting explicit cubic and contraction features instead predicts
$18$ min (native) and $76$ min (384 bits) for the negative sector at level 15;
with the same form-cache assumptions its positive-sector estimates are
$12$--$18$ min and $61$--$76$ min.  These alternatives expose model sensitivity,
especially to matrix and Gram-construction growth.  They do not form a
statistical confidence interval.  The full-feature fit has per-monomial RMS
residuals of $0.028$ s (native) and $0.068$ s (384 bits) on the saved level-10
data, but correlated features and uncounted internal Ward recursion prevent
interpreting its fitted contributions as measured stage timings.

There is no saved 136-bit PBW calibration.  In particular, the 384-bit figures
above are a comparison with a higher-precision implementation, not predictions
for a 40-digit PBW run.  Interpolating between native complex128 and FLINT by
bit count would mix different arithmetic backends and is not justified.
Nor do these estimates describe a hypothetical C++ PBW implementation.
Thus the fresh 40-digit C++ measurements establish that implementation's
runtime; the table supplies the requested approximate PBW comparison with
its precision and modelling limits stated explicitly.

\begin{samepage}
The reproducible integer counts, both fitted models, sector adjustments,
source hashes, and saved-log provenance are in
\path{C++/results/current_algorithm_2026-09-11/pbw_estimates.json}.
Regenerate them, without evaluating any block, using
\begin{verbatim}
python3 'C++/results/current_algorithm_2026-09-11/estimate_pbw.py'
\end{verbatim}
\end{samepage}
